\section{Introducción}

WhatsApp se ha convertido en un canal natural de venta para muchos negocios pequeños y medianos. Para una PyME, contestar por chat es cercano, rápido y barato; sin embargo, esa misma flexibilidad produce un problema operativo. Un cliente puede pedir una torta, enviar una foto de referencia, cambiar la fecha, olvidar la dirección y confirmar detalles en mensajes separados. El negocio entiende la intención, pero debe reconstruir manualmente el pedido, copiar datos, hacer preguntas repetidas y asumir el costo si algo se pierde en la conversación.

DeliverAI parte de esa realidad. El objetivo no es reemplazar al equipo comercial ni prometer una automatización total. El objetivo es convertir conversaciones desordenadas en pedidos estructurados, confirmados por el cliente y registrados en un sistema operacional. En términos del curso, el proyecto integra diagnóstico del problema, diseño de agente, arquitectura de integración, caso de negocio, gobierno y sustentación ejecutiva. Esa integración es importante porque un agente de IA no es solo un modelo: es un sistema socio-técnico que usa datos, herramientas, reglas, humanos y controles.

La decisión que defiende este artículo es avanzar a un \textbf{piloto controlado}. No se recomienda pasar directamente a producción. El MVP ya demuestra un flujo técnico de punta a punta, pero todavía necesita validación con clientes reales, línea base de métricas, reglas de operación, autenticación, trazabilidad completa y supervisión humana formal. La pregunta ejecutiva no es si la tecnología es interesante, sino si existe suficiente evidencia para probarla de manera acotada y gobernada.
